% Final answer for Challenge 48: Lattice Gaussian sum

For positive integer \(n\),
\[
K
=
I_{n-1}-\frac{1}{n}\mathbf{1}_{n-1}.
\]
Since the all-ones matrix has eigenvalues \(n-1\) and \(0\),
\[
\det K=\frac{1}{n},
\qquad
K^{-1}=I_{n-1}+\mathbf{1}_{n-1}.
\]

Applying \((n-1)\)-dimensional Poisson summation gives
\[
Z(n,\eta)
=
\frac{1}{\eta^{(n-1)/2}\sqrt{\det K}}
\sum_{\vec m\in\mathbb Z^{n-1}}
\exp\!\left(
-\frac{\pi}{\eta}\vec m^{\,T}K^{-1}\vec m
\right).
\]
Hence
\[
Z(n,\eta)
=
\sqrt n\,\eta^{-(n-1)/2}
\sum_{\vec m\in\mathbb Z^{n-1}}
\exp\!\left[
-\frac{\pi}{\eta}
\left(
\sum_{i=1}^{n-1}m_i^2
+
\left(\sum_{i=1}^{n-1}m_i\right)^2
\right)
\right].
\]

Introduce
\[
m_n=-\sum_{i=1}^{n-1}m_i.
\]
This is a bijective reindexing of \(\mathbb Z^{n-1}\) by integer
\(n\)-tuples satisfying \(\sum_{i=1}^n m_i=0\). Therefore
\[
Z(n,\eta)
=
\sqrt n\,\eta^{-(n-1)/2}
\sum_{\substack{\vec m\in\mathbb Z^n\\ \sum_i m_i=0}}
e^{-(\pi/\eta)\sum_i m_i^2}.
\]

Define
\[
\vartheta_\eta(t)
=
\sum_{k\in\mathbb Z}
e^{-\pi k^2/\eta}e^{2\pi i k t}.
\]
Expanding \(\vartheta_\eta(t)^n\) and integrating over one period extracts
the zero Fourier coefficient:
\[
\int_0^1\vartheta_\eta(t)^n\,dt
=
\sum_{\substack{\vec m\in\mathbb Z^n\\ \sum_i m_i=0}}
e^{-(\pi/\eta)\sum_i m_i^2}.
\]
Thus, for every positive integer \(n\),
\[
Z(n,\eta)
=
\sqrt n\,\eta^{-(n-1)/2}
\int_0^1\vartheta_\eta(t)^n\,dt.
\]

Poisson summation in one dimension also gives
\[
\vartheta_\eta(t)
=
\sqrt{\eta}
\sum_{\ell\in\mathbb Z}
e^{-\pi\eta(t-\ell)^2}
>0,
\]
so the real logarithm is well defined. By the continuation prescription
specified in the problem,
\[
Z_{\mathrm{can}}(n,\eta)
=
\sqrt n\,\eta^{-(n-1)/2}
\int_0^1
\exp\!\left[n\ln\vartheta_\eta(t)\right]\,dt.
\]

Since the constant Fourier coefficient of \(\vartheta_\eta\) is \(1\),
\[
\int_0^1\vartheta_\eta(t)\,dt=1,
\qquad
Z_{\mathrm{can}}(1,\eta)=1.
\]
Differentiating at \(n=1\) gives
\[
\left.
\frac{\partial Z_{\mathrm{can}}}{\partial n}
\right|_{n=1}
=
\frac12-\frac12\ln\eta
+
\int_0^1
\vartheta_\eta(t)\ln\vartheta_\eta(t)\,dt.
\]
Therefore
\[
F(\eta)
=
\int_0^1
\vartheta_\eta(t)\ln\vartheta_\eta(t)\,dt.
\]

At
\[
\eta=\frac{10\pi}{3},
\]
numerical evaluation gives
\[
\boxed{
F\!\left(\frac{10\pi}{3}\right)
=
0.6744991200605346131\ldots
}
\]
and hence, to eight digits after the decimal point,
\[
\boxed{F=0.67449912}.
\]
